% Luc Destombe --- licence CC BY-NC-SA 4.0
% https://creativecommons.org/licenses/by-nc-sa/4.0/deed.fr
% Fichier autonome : compiler avec pdflatex, deux fois (signets, nombre de pages).
\def\lyceeversion{1.7}\def\lyceerelectures{14}\def\lyceerodages{8}\def\lyceesceauannees{2026-2027}
\documentclass[11pt,a4paper]{article}
\makeatletter
\newif\ifeleve
\elevefalse

\newif\iflycee@palatino
\lycee@palatinotrue

\RequirePackage[utf8]{inputenc}
\RequirePackage[T1]{fontenc}
\RequirePackage[french]{babel}
\RequirePackage{etoolbox}

\newcommand{\ShorthandsOff}{\@ifundefined{shorthandoff}{}{\shorthandoff{;:!?}}}
\newcommand{\ShorthandsOn}{\@ifundefined{shorthandon}{}{\shorthandon{;:!?}}}
\ShorthandsOff
\AtBeginDocument{\ShorthandsOn}
\AtBeginEnvironment{tikzpicture}{\ShorthandsOff}

\RequirePackage{geometry}
\RequirePackage{fancyhdr}
\RequirePackage{lastpage}
\RequirePackage{multicol}
\RequirePackage{array}
\RequirePackage{enumitem}
\RequirePackage{xcolor}
\RequirePackage{needspace}

\RequirePackage{amsmath,amssymb}
\RequirePackage{mathpazo}
\renewcommand{\ttdefault}{lmtt}
\DeclareMathSizes{10.95}{10.95}{8}{5.5}
\begingroup\catcode`\_=\active \catcode`\^=\active
\gdef_#1{\sb{\scriptscriptstyle #1}}
\gdef^#1{\sp{\scriptscriptstyle\lycee@moinsepais #1}}
\endgroup
\newcommand\lycee@moins{\mathbin{%
  \setbox\z@\hbox{$\m@th\scriptscriptstyle\lycee@moinsorig$}%
  \dimen@=.55\wd\z@ \advance\dimen@-.5pt
  \hbox to.75\wd\z@{\hss$\m@th\scriptscriptstyle\vcenter{\hbox{%
    \tikz\draw[line width=.5pt,line cap=round](0,0)--(\the\dimen@,0);}}$\hss}}}
\begingroup\lccode`\~=`\- \lowercase{\endgroup\let~}\lycee@moins
\newcommand\lycee@moinsepais{\mathcode`\-="8000 }
\AtBeginDocument{\mathchardef\lycee@moinsorig=\mathcode`\-\relax}
\AtBeginDocument{\catcode`\_=12 \mathcode`\_="8000
  \catcode`\^=12 \mathcode`\^="8000 }
\newcommand\lycee@exposants{%
  \fontdimen13\textfont2=.48\fontdimen6\textfont2
  \fontdimen14\textfont2=.48\fontdimen6\textfont2
  \fontdimen15\textfont2=.48\fontdimen6\textfont2
  \fontdimen13\scriptfont2=.48\fontdimen6\scriptfont2
  \fontdimen14\scriptfont2=.48\fontdimen6\scriptfont2
  \fontdimen15\scriptfont2=.48\fontdimen6\scriptfont2}
\every@math@size\expandafter{\the\every@math@size\lycee@exposants}
\newlength{\retraitcalculs}
\setlength{\retraitcalculs}{2em}

\RequirePackage{tikz}
\usetikzlibrary{arrows.meta,calc,babel,shapes.misc}
\RequirePackage[most]{tcolorbox}
\tcbuselibrary{skins,breakable}

\definecolor{coulDef}{HTML}{1F4E79}
\definecolor{coulProp}{HTML}{7030A0}
\definecolor{coulRem}{HTML}{C55A11}
\definecolor{coulEx}{HTML}{2E7D32}
\definecolor{coulExo}{HTML}{B03A2E}
\definecolor{coulPart}{HTML}{1F4E79}
\definecolor{coulMeth}{HTML}{1B6E4F}
\definecolor{coulCourbe}{HTML}{0B5ED7}
\definecolor{coulCourbeB}{HTML}{C00000}
\definecolor{coulCourbeC}{HTML}{2E7D32}
\definecolor{coulGrille}{HTML}{8C8C8C}
\definecolor{coulRep}{HTML}{1B6E4F}
\definecolor{coulBar}{HTML}{2E7D32}

\newcommand{\R}{\mathbb{R}}
\newcommand{\Rs}{\mathbb{R}^{*}}
\renewcommand{\le}{\leqslant}
\renewcommand{\ge}{\geqslant}
\renewcommand{\approx}{\simeq}
\newcommand{\vect}[1]{\overrightarrow{#1}}
\newcommand{\norme}[1]{\left\lVert #1 \right\rVert}
\newcommand{\intervalle}[2]{\left[#1\,;#2\right]}
\RequirePackage{cancel}
\renewcommand{\CancelColor}{\color{coulExo}}
\newcommand{\fois}[1]{\times{\color{coulExo}#1}}
\newcommand{\barre}[1]{\cancel{#1}}

\tikzset{grille/.style={coulGrille,line width=0.4pt}}
\newcommand{\quadrillage}[4]{\draw[grille,step=1] (#1,#2) grid (#3,#4);}
\tikzset{
  croix/.style={cross out,draw,line width=0.6pt,inner sep=0pt,minimum size=4pt},
  croixdroite/.style={croix,rotate=45,line width=0.8pt},
}
\newcommand{\pt}[3]{\path (#1) node[croix]{} node[#2,font=\small,inner sep=2.5pt]{$#3$};}

\newcommand{\lycee@repere}[4]{%
  \draw[grille,step=1] (#1,#2) grid (#3,#4);
  \draw[-{Stealth[length=2mm]},thick] (#1,0)--({#3+0.4},0) node[below left=-1pt]{$x$};
  \draw[-{Stealth[length=2mm]},thick] (0,#2)--(0,{#4+0.4}) node[below left=-1pt]{$y$};
  \node[below left,font=\scriptsize,inner sep=1.5pt] at (0,0) {$0$};}
\newcommand{\lycee@gradx}[1]{\foreach \k/\l in {#1}{%
  \draw (\k,0.07)--(\k,-0.07) node[below,font=\scriptsize,inner sep=1.5pt]{$\l$};}}
\newcommand{\lycee@grady}[1]{\foreach \k/\l in {#1}{%
  \draw (0.07,\k)--(-0.07,\k) node[left,font=\scriptsize,inner sep=1.5pt]{$\l$};}}
\newcommand{\lycee@intff}[2]{\left[#1\,;#2\right]}
\newcommand{\lycee@intfo}[2]{\left[#1\,;#2\right[}
\newcommand{\lycee@intof}[2]{\left]#1\,;#2\right]}
\newcommand{\lycee@intoo}[2]{\left]#1\,;#2\right[}
\newcommand{\lycee@ens}[1]{\left\{#1\right\}}
\AtBeginDocument{%
  \providecommand{\repere}{\lycee@repere}%
  \providecommand{\gradx}{\lycee@gradx}%
  \providecommand{\grady}{\lycee@grady}%
  \providecommand{\Cf}{\mathcal{C}\sb{\scriptscriptstyle f}}%
  \providecommand{\Cg}{\mathcal{C}\sb{\scriptscriptstyle g}}%
  \providecommand{\intff}{\lycee@intff}%
  \providecommand{\intfo}{\lycee@intfo}%
  \providecommand{\intof}{\lycee@intof}%
  \providecommand{\intoo}{\lycee@intoo}%
  \providecommand{\ens}{\lycee@ens}}

\newlist{questions}{enumerate}{3}
\setlist[questions,1]{label=\arabic*\textdegree),leftmargin=*,itemsep=2pt,topsep=3pt}
\setlist[questions,2]{label=\alph*),leftmargin=*,itemsep=1pt,topsep=2pt}
\setlist[questions,3]{label=\roman*),leftmargin=*,itemsep=1pt,topsep=1pt}
\setlist[itemize]{label=\textbullet,leftmargin=*,itemsep=2pt,topsep=3pt}

\newcommand{\besoinplace}[1]{\par\penalty\@M\begingroup
  \dimen@=#1\relax\dimen@ii=\pagegoal\advance\dimen@ii-\pagetotal
  \ifdim\dimen@>\dimen@ii\ifdim\dimen@ii>\z@\vfil\fi\break\fi\endgroup}

\newcommand{\lycee@pied}{\thepage/\pageref{LastPage}}
\newcommand{\entete}[2]{%
  \pagestyle{fancy}\fancyhf{}%
  \fancyhead[L]{\footnotesize\color{coulPart}\textbf{#1}}%
  \fancyhead[R]{\footnotesize\color{coulPart}\textbf{#2}}%
  \fancyfoot[C]{\footnotesize\lycee@pied}%
  \renewcommand{\headrulewidth}{0.4pt}%
  \renewcommand{\headrule}{\hbox to\headwidth{%
    \color{coulPart!40}\leaders\hrule height \headrulewidth\hfill}}}

\RequirePackage{ccicons}
\newcommand{\lycee@licence}{{\scriptsize\color{gray}%
  \copyright\ Luc Destombe --- \ccbyncsa\ CC BY-NC-SA 4.0}}
\newif\iflycee@licence \lycee@licencetrue
\newcommand{\sanslicence}{\lycee@licencefalse}
\AtBeginDocument{\iflycee@licence\AddToHookNext{shipout/foreground}{%
  \put(0,0){\rlap{\kern\dimexpr1in+\hoffset+\oddsidemargin+\textwidth\relax
    \llap{\raisebox{-\dimexpr1in+\voffset+\topmargin+\headheight+\headsep
      +\textheight+\footskip\relax}{\lycee@licence}}}}}\fi}

\newcommand{\lycee@site}{faire-des-maths.fr}
\newcommand{\lycee@version}{{\scriptsize\color{gray}\lycee@site
  \ifdefined\lyceeversion\ --- Version~\lyceeversion\fi}}
\AtBeginDocument{\iflycee@licence\AddToHookNext{shipout/foreground}{%
  \put(0,0){\rlap{\kern\dimexpr1in+\hoffset+\oddsidemargin\relax
    \raisebox{-\dimexpr1in+\voffset+\topmargin+\headheight+\headsep
      +\textheight+\footskip\relax}{\lycee@version}}}}\fi}

\newcommand{\formulesserrees}{%
  \setlength{\abovedisplayskip}{3pt}\setlength{\belowdisplayskip}{3pt}%
  \setlength{\abovedisplayshortskip}{2pt}\setlength{\belowdisplayshortskip}{3pt}}

\setlength{\parindent}{0pt}

\RequirePackage[implicit=false,hidelinks,pdfencoding=auto,psdextra,bookmarksopen,
  bookmarksopenlevel=1,pdfcreator={},pdfproducer={}]{hyperref}
\RequirePackage{bookmark}
\hypersetup{pdfauthor={Luc Destombe},pdfkeywords={CC BY-NC-SA 4.0}}
\newcounter{lycee@signet}
\newcommand{\signet}[3][]{\stepcounter{lycee@signet}%
  \edef\lycee@dest{\if\relax\detokenize{#1}\relax
    signet.\arabic{lycee@signet}\else#1\fi}%
  \hypertarget{\lycee@dest}{}\bookmark[level=#2,dest=\lycee@dest]{#3}}
\pdfstringdefDisableCommands{%
  \def\R{R}\def\vect#1{#1}\def\Cf{Cf}\def\Cg{Cg}\def\fois#1{×#1}%
  \def\barre#1{#1}\def\intervalle#1#2{[#1;#2]}\def\intff#1#2{[#1;#2]}%
  \def\textsuperscript#1{#1}\def\dfrac#1#2{#1/#2}\def\frac#1#2{#1/#2}%
  \def\vec#1{#1⃗}}%
\newcommand{\res}[1]{#1}
\geometry{top=2cm,bottom=1cm,left=1cm,right=1cm,headheight=14pt,footskip=0.6cm}

\RequirePackage{pgfplots}
\pgfplotsset{compat=1.18}
\RequirePackage{tkz-tab}

\newcounter{partiecours}
\renewcommand{\thepartiecours}{\Roman{partiecours}}
\newcommand{\partiecours}[1]{%
  \refstepcounter{partiecours}%
  \Needspace*{8\baselineskip}%
  \bigskip
  \noindent\begin{tcolorbox}[enhanced,colback=coulPart,colframe=coulPart,
      boxrule=0pt,arc=2pt,left=8pt,right=8pt,top=4pt,bottom=4pt,
      before skip=6pt,after skip=10pt]
    \color{white}\bfseries\large \thepartiecours{} --- #1%
    \signet[partie-\arabic{partiecours}]{1}{\thepartiecours{} --- #1}
  \end{tcolorbox}%
  \addcontentsline{toc}{section}{\thepartiecours{} --- #1}%
}

\newtcolorbox{boitecours}[3][]{%
  enhanced,breakable,
  colback=white,colframe=#2,coltitle=white,
  fonttitle=\bfseries,
  attach boxed title to top left={xshift=8pt,yshift=-\tcboxedtitleheight/2},
  boxed title style={colback=#2,arc=2pt,boxrule=0pt},
  boxrule=0.9pt,arc=2pt,
  left=8pt,right=8pt,top=8pt,bottom=6pt,
  title={#3},#1}

\newcounter{methode}
\newenvironment{definitionbox}[1][Définition]%
  {\begin{boitecours}[phantom={\signet{2}{#1}}]{coulDef}{#1}}{\end{boitecours}}
\newenvironment{proprietebox}[1][Propriété]%
  {\begin{boitecours}[phantom={\signet{2}{#1}}]{coulProp}{#1}}{\end{boitecours}}
\newenvironment{methodebox}[1][Méthode]{\stepcounter{methode}%
  \begin{boitecours}[phantom={\signet[methode-\arabic{methode}]{2}{#1}}]{coulMeth}{#1}}%
  {\end{boitecours}}

\newcommand{\etape}[1]{\par\smallskip\textbf{\color{coulMeth}#1}\par\nobreak\smallskip}
\newcommand{\enonce}[1]{\emph{#1}\par\smallskip}
\newcommand{\reponse}[1]{\par\smallskip
  {\footnotesize\itshape\color{black!55}Réponses : #1}}
\newcommand{\demo}[1]{\textbf{\color{coulMeth}Démonstration.} #1}
\newcommand{\afaire}[1]{%
  \par\smallskip
  \noindent{\small\color{coulExo}$\blacktriangleright$\ \textbf{À faire :}
    fiche d'exercices du chapitre, #1.}\par\smallskip}

\newenvironment{calculs}{\setlength{\jot}{5pt}\@fleqntrue\@mathmargin\retraitcalculs
  \setlength{\abovedisplayskip}{4pt}\setlength{\belowdisplayskip}{4pt}%
  \setlength{\abovedisplayshortskip}{4pt}\setlength{\belowdisplayshortskip}{4pt}%
  \csname alignat*\endcsname{2}}%
  {\csname endalignat*\endcsname}
\newcommand{\rem}[1]{\text{\small\itshape\color{black!60}#1}}
\newcommand{\explique}[2]{\par\smallskip\noindent
  \begin{minipage}[t]{0.57\linewidth}\raggedright #1\end{minipage}\hfill
  \begin{minipage}[t]{0.40\linewidth}\raggedright\leavevmode
    {\small\itshape\color{black!60}#2}\end{minipage}\par}
\newcommand{\cas}[1]{\par\smallskip\noindent\textbf{\color{coulExo}#1}\vadjust{\nobreak}}
\newcommand{\calc}[1]{\par\smallskip\textbf{\color{coulExo}#1}\par\nobreak}

\newtcolorbox{remarquebox}[1][Remarque]{%
  enhanced,breakable,colback=white,colframe=coulRem,
  boxrule=0pt,leftrule=2.5pt,arc=0pt,outer arc=0pt,
  left=8pt,right=6pt,top=5pt,bottom=5pt,
  before upper={\textbf{\color{coulRem}#1}\par\smallskip}}

\newtcolorbox{exemplebox}[1][Exemple]{%
  enhanced,breakable,colback=white,colframe=coulEx,
  boxrule=0pt,leftrule=2.5pt,arc=0pt,outer arc=0pt,
  left=8pt,right=6pt,top=5pt,bottom=5pt,
  before upper={\textbf{\color{coulEx}#1}\par\smallskip}}

\pgfplotsset{
  repere/.style={
    axis lines=middle,
    axis line style={-{Stealth[length=2.4mm]},thick},
    every axis x label/.style={at={(ticklabel* cs:1.0)},anchor=west},
    every axis y label/.style={at={(ticklabel* cs:1.0)},anchor=south},
    label style={font=\footnotesize},
    tick label style={font=\scriptsize},
    grid=both,
    major grid style={grille},
    minor tick num=0,
    tick align=inside,
    clip mode=individual,
  },
}
\tikzset{
  courbe/.style={coulCourbe,line width=1.1pt,smooth},
  courbeB/.style={coulCourbeB,line width=1.1pt,smooth},
  courbeC/.style={coulCourbeC,line width=1.1pt,smooth},
}

\newlist{etapes}{enumerate}{1}
\setlist[etapes]{label=\textbf{\arabic*.},leftmargin=*,itemsep=1pt,topsep=2pt}

\newlist{expressions}{enumerate}{1}
\setlist[expressions]{label={\Alph*\ $=$},leftmargin=*,itemsep=3pt,topsep=3pt}

\newcounter{exo}
\renewcommand{\theexo}{\ifnum\value{exo}<10 0\fi\arabic{exo}}
\newtcolorbox{exercice}[1][]{%
  enhanced,breakable,
  colback=white,colframe=coulExo,
  boxrule=0.9pt,arc=2pt,
  left=8pt,right=8pt,top=8pt,bottom=6pt,
  coltitle=white,fonttitle=\bfseries,
  attach boxed title to top left={xshift=8pt,yshift=-\tcboxedtitleheight/2},
  boxed title style={colback=coulExo,arc=2pt,boxrule=0pt},
  phantom={\signet{2}{Exercice \arabic{exo}}},
  title={Exercice \theexo\ifx\relax#1\relax\else{} --- #1\fi}}
\newenvironment{exo}[1][\relax]{\refstepcounter{exo}\begin{exercice}[#1]}{\end{exercice}}

  \newenvironment{correction}{\par}{\par}
  \newcommand{\autableau}{}
  \newcommand{\etapeexemples}{\etape{Exemples corrigés}}
  \newcommand{\etapetableau}[1]{\etape{#1}}
  \newenvironment{exempletableau}[1][Exemple]%
    {\begin{exemplebox}[#1]}{\end{exemplebox}}

\RequirePackage{comment}
\excludecomment{correctionavj}

\newcommand{\coursEntete}{}
\newcommand{\coursNiveau}{}
\newcommand{\coursTitre}{}
\newcommand{\coursSurTitre}{}
\newcommand{\coursSousTitre}{}

\newcommand{\enteteCours}[1]{\renewcommand{\coursEntete}{#1}}
\newcommand{\chapitrecours}[2]{\enteteCours{Chapitre #1 --- #2}}
\newcommand{\niveaucours}[1]{\renewcommand{\coursNiveau}{#1}}
\newcommand{\titrecours}[3][]{%
  \renewcommand{\coursTitre}{#2}%
  \renewcommand{\coursSurTitre}{#1}%
  \renewcommand{\coursSousTitre}{#3}}

\pagestyle{fancy}
\fancyhf{}
\fancyhead[L]{\footnotesize\color{coulPart}\textbf{\coursEntete}}
\fancyhead[R]{\footnotesize\color{coulPart}\textbf{\coursNiveau}}
\fancyfoot[C]{\footnotesize\thepage}
\renewcommand{\headrulewidth}{0.4pt}
\renewcommand{\headrule}{\hbox to\headwidth{\color{coulPart!40}\leaders\hrule height \headrulewidth\hfill}}

\setlength{\parskip}{2pt}

\AtBeginDocument{%
  \begin{center}
    \begin{tcolorbox}[enhanced,width=0.72\linewidth,colback=white,colframe=coulPart,
        boxrule=1.6pt,arc=3pt,halign=center,top=10pt,bottom=10pt]
      \ifx\coursSurTitre\empty
        {\Huge\bfseries\color{coulPart} \coursTitre}\\[4pt]
      \else
        {\Huge\bfseries\color{coulPart} \coursTitre}\\[3pt]
        {\Large\bfseries\color{coulPart} \coursSurTitre}\\[5pt]
      \fi
      {\normalsize \coursSousTitre}%
      
    \end{tcolorbox}
  \end{center}%
}

\usetikzlibrary{decorations.text}
\definecolor{lycee@metal}{HTML}{1F4E79}
\definecolor{lycee@manche}{HTML}{8A5A2B}
\definecolor{lycee@etincelle}{HTML}{F28C28}
\newcommand{\lycee@pluriel}[2]{#1 #2\ifnum#1>1 s\fi}
\newcommand{\lycee@sceau}{%
  \edef\lycee@bas{\lycee@pluriel{\lyceerelectures}{relecture} et
    \lycee@pluriel{\lyceerodages}{rodage}}%
  \begin{tikzpicture}[line join=round]
    \fill[coulPart!10,draw=coulPart,line width=.8pt] (0:2.2)
      \foreach \i in {1,...,71}{ -- (\i*5:{isodd(\i)?2.1:2.2}) } -- cycle;
    \draw[coulPart,line width=1.2pt] (0,0) circle (1.93);
    \draw[coulPart,line width=.4pt] (0,0) circle (1.85);
    \draw[coulPart,line width=.4pt] (0,0) circle (1.2);
    \path[decorate,decoration={text along path,text align=center,
        text={|\fontsize{9.5}{11}\selectfont\bfseries\color{coulPart}|TEST{\'E} AVEC MES {\'E}L{\`E}VES}}]
      (180:1.45) arc (180:0:1.45);
    \path[decorate,decoration={text along path,text align=center,
        text/.expanded={|\noexpand\fontsize{8.5}{10}\noexpand\selectfont\noexpand\color{coulPart}|\lycee@bas}}]
      (195:1.68) arc (195:345:1.68);
    \begin{scope}[x=.25mm,y=-.25mm]
      \begin{scope}[shift={(-32,-40)}]%
        \fill[lycee@metal] (23,31) rectangle (41,49);
        \draw[lycee@etincelle,line width=.5pt,line cap=round]
          (20,30)--(15,27) (19,34)--(14,35) (44,30)--(49,27) (45,34)--(50,35);
        \fill[lycee@manche,rounded corners=.3mm] (38,24.6) rectangle (60,27.6);
        \fill[lycee@metal,rounded corners=.4mm] (25,20) rectangle (39,31);
        \fill[lycee@etincelle,rounded corners=.2mm] (25,20) rectangle (28,31);
      \end{scope}
    \end{scope}
    \ifdefined\lyceesceauplusieurs\def\lycee@annee{Ann{\'e}es~:}\else
      \def\lycee@annee{Ann{\'e}e}\fi
    \path[decorate,decoration={text along path,text align=center,
        text/.expanded={|\noexpand\fontsize{8}{10}\noexpand\selectfont\noexpand\color{coulPart}|%
          \unexpanded\expandafter{\lycee@annee} \unexpanded\expandafter{\lyceesceauannees}}}]
      (95:1.08) arc (95:445:1.08);
  \end{tikzpicture}}
\AtEndDocument{\ifdefined\lyceesceauannees
  \par\besoinplace{5cm}\vspace{1.5em}%
  \begin{center}\lycee@sceau\end{center}\fi}
\makeatother

\usepackage{xlop}
\opset{decimalsepsymbol={,}}
\newcommand{\divpose}[3]{\begin{tikzpicture}
  \node{\opdiv[maxdivstep=#1,safedivstep=false,displayintermediary=all]{#2}{#3}};
  \end{tikzpicture}}
\newcommand{\divposesuite}[3]{\begin{tikzpicture}
  \node[inner sep=0pt](d){\opdiv[maxdivstep=#1,safedivstep=false,displayintermediary=all]{#2}{#3}};
  \node[anchor=base west,inner sep=1pt] at ([yshift=-1.8\baselineskip]d.north east) {\ldots};
  \pgfmathsetmacro{\colonne}{#1+(#2>9?2:1)-0.5}
  \node[anchor=base,inner sep=0pt] at ([xshift=\colonne*2ex,yshift=-\baselineskip]d.south west) {\ldots};
  \end{tikzpicture}}
\newenvironment{division}[1][0.62]{\par\smallskip\noindent\begin{minipage}[t]{#1\linewidth}\vspace{0pt}\raggedright}%
  {\end{minipage}\hfill}
\newenvironment{figdivision}[1][0.34]{\begin{minipage}[t]{#1\linewidth}\vspace{0pt}\centering}%
  {\end{minipage}\par\smallskip}

\chapitrecours{1}{Fractions et puissances}
\niveaucours{Seconde}
\titrecours{FRACTIONS ET PUISSANCES}{Cours --- Seconde}

\begin{document}

\begin{remarquebox}[Comment utiliser ce cours]
Chaque partie commence par ce qu'il faut \textbf{savoir} (définitions, propriétés), puis
vient ce qu'il faut \textbf{savoir faire} : une méthode, avec des
\textbf{exemples corrigés} faits en classe, puis des exercices
\og\textbf{À vous de jouer}\fg{} à chercher seul.

\end{remarquebox}

\partiecours{Développement décimal limité et illimité d'une fraction}

\begin{definitionbox}[Définition 1 --- Développement décimal]
On obtient l'écriture décimale d'une fraction $\dfrac{a}{b}$ en posant la division de $a$
par $b$ : c'est son \textbf{développement décimal}.
\begin{itemize}
  \item Si un reste devient nul, la division se termine : le développement est
        \textbf{limité} et donne un \textbf{nombre rationnel décimal}.
        \begin{division}
        Exemple : on pose la division de $7$ par $8$.

        \smallskip
        On obtient $\dfrac{7}{8}=0{,}875=\dfrac{875}{1000}$ : cette fraction est égale à
        une \textbf{fraction décimale} (de dénominateur $10$, $100$, $1000$\ldots).
        \end{division}
        \begin{figdivision}
        \divpose{4}{7}{8}
        \end{figdivision}
  \item Sinon, la division ne se termine pas : le développement est \textbf{illimité}. Un
        reste revient forcément (car il n'y a que $b-1$ restes non nuls possibles).
        \begin{division}
        Exemple : on pose la division de $1$ par $3$.

        \smallskip
        On obtient $\dfrac{1}{3}=0{,}333\ldots$ Le bloc de chiffres qui se répète sans
        fin est appelé la \textbf{période} : ici, $\dfrac{1}{3}$ a pour période $3$.

        \smallskip
        La fraction n'est pas égale à une fraction décimale : elle donne un
        \textbf{nombre rationnel non décimal}.
        \end{division}
        \begin{figdivision}
        \divposesuite{4}{1}{3}
        \end{figdivision}
\end{itemize}
\end{definitionbox}

\begin{methodebox}[Méthode 1 --- Écrire le développement décimal d'une fraction]
On pose la division comme à l'école : on écrit la virgule au quotient, puis, à chaque
étape, on abaisse un $0$ à côté du reste. Si un reste vaut $0$, la division se termine. Si
un reste déjà obtenu revient, elle ne se termine pas : les mêmes chiffres reviennent sans
fin.
\etapeexemples
Poser la division, puis donner le développement décimal de chaque fraction ; dire s'il
est limité ou illimité et, s'il est illimité, donner sa période. Conclure : nombre
rationnel décimal ou nombre rationnel non décimal ?
\quad a)~$\dfrac{7}{8}$ \qquad b)~$\dfrac{2}{3}$ \qquad c)~$\dfrac{13}{6}$ \qquad
d)~$\dfrac{3}{7}$
\begin{correction}
\cas{a)}
\begin{division}
Le reste $0$ apparaît : la division se termine.
\begin{calculs}
\dfrac{7}{8} &= 0{,}875 &\quad& \rem{développement limité}
\end{calculs}
\par\smallskip
Donc $\dfrac{7}{8}$ est un nombre rationnel décimal.
\end{division}
\begin{figdivision}
\divpose{4}{7}{8}
\end{figdivision}
\cas{b)}
\begin{division}
Le reste $2$ revient : la division ne se termine pas, le chiffre $6$ revient sans fin.
\begin{calculs}
\dfrac{2}{3} &= 0{,}666\ldots &\quad& \rem{illimité, de période $6$}
\end{calculs}
\par\smallskip
Donc $\dfrac{2}{3}$ est un nombre rationnel non décimal.
\end{division}
\begin{figdivision}
\divposesuite{4}{2}{3}
\end{figdivision}
\cas{c)}
\begin{division}
Le reste $4$ revient : la division ne se termine pas, le chiffre $6$ revient sans fin.
\begin{calculs}
\dfrac{13}{6} &= 2{,}1666\ldots &\quad& \rem{illimité, de période $6$}
\end{calculs}
\explique{Le chiffre $1$ ne revient pas.}{la période peut commencer plus loin}
\par\smallskip
Donc $\dfrac{13}{6}$ est un nombre rationnel non décimal.
\end{division}
\begin{figdivision}
\divposesuite{5}{13}{6}
\end{figdivision}
\cas{d)}
\begin{division}[0.5]
Après six restes différents, le reste $3$ du départ revient : la division ne se termine pas, on repart sur~$42\ldots$
\begin{calculs}
\dfrac{3}{7} &= 0{,}428571\,42\ldots &\quad& \rem{illimité, de période $428571$}
\end{calculs}
\par\smallskip
Donc $\dfrac{3}{7}$ est un nombre rationnel non décimal.
\end{division}
\begin{figdivision}[0.47]
\divposesuite{9}{3}{7}
\end{figdivision}
\end{correction}

\etape{À vous de jouer}
Mêmes questions pour :
\quad e)~$\dfrac{9}{4}$ \qquad f)~$\dfrac{5}{9}$ \qquad g)~$\dfrac{7}{6}$ \qquad
h)~$\dfrac{2}{7}$
\reponse{e) $2{,}25$, limité, nombre rationnel décimal ;\; f)~$0{,}555\ldots$, période $5$,
nombre rationnel non décimal ;\; g)~$1{,}1666\ldots$, période $6$, nombre rationnel non
décimal ;\; h)~$0{,}285714\,285714\ldots$, période $285714$, nombre rationnel non décimal.}
\end{methodebox}

\begin{definitionbox}[Définition 2 --- Valeur approchée par défaut, par excès]
Le développement décimal de $\dfrac{3}{7}$ commence par $0{,}428\ldots$ : on a
$0{,}42<\dfrac{3}{7}<0{,}43$. On dit que $0{,}42$ est la \textbf{valeur approchée par
défaut} de $\dfrac{3}{7}$ au centième, et que $0{,}43$ est sa \textbf{valeur approchée par
excès} au centième.
\end{definitionbox}

\begin{methodebox}[Méthode 2 --- Donner une valeur approchée d'une fraction]
On écrit le développement décimal (méthode 1) avec un chiffre de plus que la précision
demandée ; on coupe après le chiffre demandé (par défaut), puis on ajoute $1$ à ce
chiffre (par excès).
\etapeexemples
Donner la valeur approchée par défaut et la valeur approchée par excès, puis encadrer la
fraction :
\quad a)~$\dfrac{3}{7}$ au centième \qquad b)~$\dfrac{13}{6}$ au millième \qquad
c)~$\dfrac{2}{3}$ au dixième
\begin{correction}
\cas{a)}
\begin{calculs}
\dfrac{3}{7} &= 0{,}428\ldots &\qquad& \rem{méthode 1 ; deux chiffres après la virgule, plus un}\\
\text{par défaut} &: 0{,}42 &\qquad& \rem{on coupe après le centième}\\
\text{par excès} &: 0{,}43 &\qquad& \rem{on ajoute $1$ au chiffre des centièmes}\\
0{,}42 &< \dfrac{3}{7}<0{,}43 &&
\end{calculs}
\cas{b)}
\begin{calculs}
\dfrac{13}{6} &= 2{,}1666\ldots &\qquad& \rem{méthode 1}\\
\text{par défaut} &: 2{,}166 &\qquad& \rem{on coupe après le millième}\\
\text{par excès} &: 2{,}167 &\qquad& \rem{on ajoute $1$ au chiffre des millièmes}\\
2{,}166 &< \dfrac{13}{6}<2{,}167 &&
\end{calculs}
\cas{c)}
\begin{calculs}
\dfrac{2}{3} &= 0{,}66\ldots &\qquad& \rem{méthode 1}\\
\text{par défaut} &: 0{,}6 &\qquad& \rem{on coupe après le dixième}\\
\text{par excès} &: 0{,}7 &\qquad& \rem{on ajoute $1$ au chiffre des dixièmes}\\
0{,}6 &< \dfrac{2}{3}<0{,}7 &&
\end{calculs}
\end{correction}

\etape{À vous de jouer}
Mêmes questions pour :
\quad d)~$\dfrac{2}{7}$ au centième \qquad e)~$\dfrac{7}{6}$ au millième \qquad
f)~$\dfrac{5}{9}$ au dixième
\reponse{d) $0{,}28<\frac27<0{,}29$ ;\; e)~$1{,}166<\frac{7}{6}<1{,}167$ ;\;
f)~$0{,}5<\frac{5}{9}<0{,}6$.}
\end{methodebox}

\afaire{exercices 1 à 4}

\partiecours{Simplification de fractions}

\begin{definitionbox}[Définition 3 --- Nombre premier]
Un \textbf{nombre premier} est un entier qui a exactement deux diviseurs : $1$ et
lui-même. Les premiers sont $2$, $3$, $5$, $7$, $11$, $13$, $17$, $19$, $23$, $29$\ldots

\smallskip
$1$ n'est pas premier (un seul diviseur) ; $15$ n'est pas premier ($15=3\times 5$).
\end{definitionbox}

\begin{proprietebox}[Propriété 1 --- Décomposition en produit de facteurs premiers]
Tout entier supérieur ou égal à $2$ s'écrit comme un produit de nombres premiers, et
d'une seule façon à l'ordre près : $60=2\times 2\times 3\times 5$.
\end{proprietebox}

\begin{remarquebox}[Rappel : critères de divisibilité]
Un entier est divisible par $2$ si son chiffre des unités est pair ; par $3$ si la somme
de ses chiffres est un multiple de $3$ ; par $5$ si son chiffre des unités est $0$ ou $5$.
Pour $7$, $11$, $13$\ldots, on pose la division.
\end{remarquebox}

\begin{methodebox}[Méthode 3 --- Décomposer un entier en produit de facteurs premiers]
On divise par le plus petit nombre premier possible ($2$, puis $3$, puis $5$\ldots), puis
on recommence avec le quotient, jusqu'à obtenir $1$ : ce sont les \textbf{divisions
successives}.
\etapeexemples
Décomposer en produit de facteurs premiers :
\quad a)~$60$ \qquad b)~$126$ \qquad c)~$360$ \qquad d)~$231$
\begin{correction}
\cas{a)}
\[
\renewcommand{\arraystretch}{1.3}
\begin{array}{r|l}
60 & 2 \quad \rem{$60$ est pair}\\
30 & 2 \quad \rem{$30$ est pair}\\
15 & 3 \quad \rem{$1+5=6$, multiple de $3$}\\
5 & 5 \quad \rem{$5$ est premier}\\
1 &
\end{array}
\]
\explique{$60=2\times 2\times 3\times 5$.}{on multiplie les diviseurs de la colonne de droite}
\cas{b)}
\[
\renewcommand{\arraystretch}{1.3}
\begin{array}{r|l}
126 & 2 \quad \rem{$126$ est pair}\\
63 & 3 \quad \rem{$6+3=9$, multiple de $3$}\\
21 & 3 \quad \rem{$2+1=3$}\\
7 & 7 \quad \rem{$7$ est premier}\\
1 &
\end{array}
\]
\explique{$126=2\times 3\times 3\times 7$.}{}
\cas{c)}
\[
\renewcommand{\arraystretch}{1.3}
\begin{array}{r|l}
360 & 2 \quad \rem{tant que le nombre est pair, on divise par $2$}\\
180 & 2\\
90 & 2\\
45 & 3 \quad \rem{$4+5=9$ : on passe à $3$}\\
15 & 3\\
5 & 5\\
1 &
\end{array}
\]
\explique{$360=2\times 2\times 2\times 3\times 3\times 5$.}{}
\cas{d)}
\[
\renewcommand{\arraystretch}{1.3}
\begin{array}{r|l}
231 & 3 \quad \rem{$231$ est impair ; $2+3+1=6$, multiple de $3$}\\
77 & 7 \quad \rem{ni $3$ ni $5$ : on essaie $7$, $77=7\times 11$}\\
11 & 11 \quad \rem{$11$ est premier}\\
1 &
\end{array}
\]
\explique{$231=3\times 7\times 11$.}{}
\end{correction}

\etape{À vous de jouer}
Décomposer en produit de facteurs premiers :
\quad e)~$90$ \qquad f)~$150$ \qquad g)~$540$ \qquad h)~$455$
\reponse{e) $90=2\times 3\times 3\times 5$ ;\; f)~$150=2\times 3\times 5\times 5$ ;\;
g)~$540=2\times 2\times 3\times 3\times 3\times 5$ ;\; h)~$455=5\times 7\times 13$.}
\end{methodebox}

\begin{definitionbox}[Définition 4 --- PGCD, fraction irréductible]
Le \textbf{PGCD} de deux entiers est leur plus grand diviseur commun ; on le note
$\mathrm{PGCD}(60\,;126)$.

\smallskip
Une fraction est \textbf{irréductible} quand son numérateur et son dénominateur n'ont pas
d'autre diviseur commun que $1$ : on ne peut plus la simplifier.
\end{definitionbox}

\begin{proprietebox}[Propriété 2 --- Calcul du PGCD]
Le PGCD est le produit des facteurs premiers \textbf{communs} aux deux décompositions,
chacun pris autant de fois qu'il est commun. En simplifiant une fraction par le PGCD, on
obtient une fraction irréductible.
\end{proprietebox}

\begin{methodebox}[Méthode 4 --- Simplifier une fraction avec le PGCD]
\etapeexemples
Calculer le PGCD du numérateur et du dénominateur, puis écrire la fraction sous forme
irréductible :
\quad a)~$\dfrac{60}{126}$ \qquad b)~$\dfrac{84}{360}$ \qquad c)~$\dfrac{231}{385}$
\begin{correction}
\cas{a)}
\begin{calculs}
60 &= \mathbf{2}\times 2\times \mathbf{3}\times 5 &\qquad& \rem{méthode 3 ; facteurs communs en gras}\\
126 &= \mathbf{2}\times \mathbf{3}\times 3\times 7 &&\\
\mathrm{PGCD}(60\,;126) &= 2\times 3 &\qquad& \rem{facteurs communs : un $2$ et un $3$}\\
\mathrm{PGCD}(60\,;126) &= 6 &&\\
\dfrac{60}{126} &= \dfrac{\barre{2}\times 2\times\barre{3}\times 5}{\barre{2}\times\barre{3}\times 3\times 7} &\qquad& \rem{on barre les facteurs communs, en haut et en bas}\\
\dfrac{60}{126} &= \dfrac{10}{21} &\qquad& \rem{on a simplifié par le PGCD $6$ : irréductible}
\end{calculs}
\cas{b)}
\begin{calculs}
84 &= \mathbf{2}\times \mathbf{2}\times \mathbf{3}\times 7 &\qquad& \rem{divisions successives : $84$, $42$, $21$, $7$}\\
360 &= \mathbf{2}\times \mathbf{2}\times 2\times \mathbf{3}\times 3\times 5 &\qquad& \rem{méthode 3, exemple c)}\\
\mathrm{PGCD}(84\,;360) &= 2\times 2\times 3 &\qquad& \rem{deux $2$ et un $3$ sont communs}\\
\mathrm{PGCD}(84\,;360) &= 12 &&\\
\dfrac{84}{360} &= \dfrac{\barre{2}\times\barre{2}\times\barre{3}\times 7}{\barre{2}\times\barre{2}\times 2\times\barre{3}\times 3\times 5} &&\\
\dfrac{84}{360} &= \dfrac{7}{30} &\qquad& \rem{$2\times 3\times 5=30$}
\end{calculs}
\cas{c)}
\begin{calculs}
231 &= 3\times \mathbf{7}\times \mathbf{11} &\qquad& \rem{méthode 3, exemple d)}\\
385 &= 5\times \mathbf{7}\times \mathbf{11} &\qquad& \rem{divisions successives : $385$, $77$, $11$}\\
\mathrm{PGCD}(231\,;385) &= 7\times 11 &\qquad& \rem{deux facteurs communs, cachés}\\
\mathrm{PGCD}(231\,;385) &= 77 &&\\
\dfrac{231}{385} &= \dfrac{3\times\barre{7}\times\barre{11}}{5\times\barre{7}\times\barre{11}} &&\\
\dfrac{231}{385} &= \dfrac{3}{5} &&
\end{calculs}
\end{correction}

\etape{À vous de jouer}
Mêmes questions pour :
\quad d)~$\dfrac{90}{150}$ \qquad e)~$\dfrac{84}{540}$ \qquad f)~$\dfrac{154}{455}$
\reponse{d) PGCD $30$, $\frac{90}{150}=\frac35$ ;\; e) PGCD $12$, $\frac{84}{540}=\frac{7}{45}$ ;\;
f) PGCD $7$, $\frac{154}{455}=\frac{22}{65}$.}
\end{methodebox}

\afaire{exercices 5 à 8}

\partiecours{Additionner et soustraire des fractions}

\begin{proprietebox}[Propriété 3 --- Le signe d'une fraction]
Pour $b\ne 0$ : \quad $\dfrac{-a}{b}=\dfrac{a}{-b}=-\dfrac{a}{b}$ \quad et \quad
$\dfrac{-a}{-b}=\dfrac{a}{b}$.

\smallskip
On écrit le signe moins d'une fraction au numérateur ou devant le trait de fraction,
jamais au dénominateur. Un signe moins devant une fraction s'applique à \textbf{tout}
le numérateur : on le met entre parenthèses.
\end{proprietebox}

\besoinplace{5cm}
\begin{exemplebox}[Exemple]
\begin{calculs}
\dfrac{5}{-8} &= -\dfrac{5}{8} &\qquad& \rem{on ne laisse pas le moins au dénominateur}\\
\dfrac{-6}{-7} &= \dfrac{6}{7} &\qquad& \rem{deux signes moins : la fraction est positive}\\
-\dfrac{2-9}{4} &= -\dfrac{-7}{4} &\qquad& \rem{le moins porte sur tout le numérateur, calculé d'abord}\\
-\dfrac{2-9}{4} &= \dfrac{7}{4} &\qquad& \rem{deux signes moins}
\end{calculs}
\end{exemplebox}

\begin{proprietebox}[Propriété 4 --- Somme de fractions]
Pour additionner (ou soustraire) des fractions, on les écrit avec le \textbf{même
dénominateur}, puis on additionne (ou on soustrait) les numérateurs ; le dénominateur
ne change pas :
\[
  \dfrac{a}{d}+\dfrac{b}{d}=\dfrac{a+b}{d} \qquad\qquad \dfrac{a}{d}-\dfrac{b}{d}=\dfrac{a-b}{d}
\]
\end{proprietebox}

\begin{methodebox}[Méthode 5 --- Additionner et soustraire des fractions]
\etapeexemples
Calculer et donner le résultat sous forme simplifiée ; pour d), $x$ désigne un nombre,
écrire l'expression avec une seule fraction :
\begin{questions}[label=\alph*)]
  \item $A=\dfrac{-7}{-6}+\dfrac{-3}{12}$
  \item $B=\dfrac{4}{5}-\dfrac{-2}{3}$
  \item $C=\dfrac{1}{3}-\dfrac{1}{6}+\dfrac{1}{9}-\dfrac{1}{18}$
  \item $D=3-\dfrac{x-1}{5}$
\end{questions}
\begin{correction}
\cas{a)}
\begin{calculs}
A &= \dfrac{-7}{-6}+\dfrac{-3}{12} &\qquad& \rem{on recopie}\\
A &= \dfrac{7}{6}-\dfrac{3}{12} &\qquad& \rem{on simplifie les signes $-$}\\
A &= \dfrac{7\fois{2}}{6\fois{2}}-\dfrac{3}{12} &\qquad& \rem{$12$ est un multiple de $6$}\\
A &= \dfrac{14}{12}-\dfrac{3}{12} &&\\
A &= \dfrac{11}{12} &\qquad& \rem{on soustrait les numérateurs}
\end{calculs}
\cas{b)}
\begin{calculs}
B &= \dfrac{4}{5}-\dfrac{-2}{3} &\qquad& \rem{on recopie}\\
B &= \dfrac{4}{5}+\dfrac{2}{3} &\qquad& \rem{soustraire $-\frac23$, c'est ajouter $\frac23$}\\
B &= \dfrac{4\fois{3}}{5\fois{3}}+\dfrac{2\fois{5}}{3\fois{5}} &\qquad& \rem{dénominateur commun : $5\times 3=15$}\\
B &= \dfrac{12}{15}+\dfrac{10}{15} &&\\
B &= \dfrac{22}{15} &\qquad& \rem{$22$ et $15$ n'ont pas de diviseur commun}
\end{calculs}
\cas{c)}
\begin{calculs}
C &= \dfrac{1}{3}-\dfrac{1}{6}+\dfrac{1}{9}-\dfrac{1}{18} &\qquad& \rem{$18$ est un multiple de $3$, $6$ et $9$}\\
C &= \dfrac{1\fois{6}}{3\fois{6}}-\dfrac{1\fois{3}}{6\fois{3}}+\dfrac{1\fois{2}}{9\fois{2}}-\dfrac{1}{18} &\qquad& \rem{tout sur $18$}\\
C &= \dfrac{6}{18}-\dfrac{3}{18}+\dfrac{2}{18}-\dfrac{1}{18} &&\\
C &= \dfrac{4}{18} &\qquad& \rem{$6-3+2-1=4$}\\
C &= \dfrac{2}{9} &\qquad& \rem{on simplifie par $2$}
\end{calculs}
\cas{d)}
\begin{calculs}
D &= 3-\dfrac{x-1}{5} &\qquad& \rem{on recopie}\\
D &= \dfrac{3\fois{5}}{1\fois{5}}-\dfrac{x-1}{5} &\qquad& \rem{$3=\frac31$, écrit avec le dénominateur $5$}\\
D &= \dfrac{15}{5}-\dfrac{x-1}{5} &&\\
D &= \dfrac{15-(x-1)}{5} &\qquad& \rem{parenthèses : le moins porte sur tout le numérateur}\\
D &= \dfrac{15-x+1}{5} &\qquad& \rem{$-(x-1)=-x+1$}\\
D &= \dfrac{16-x}{5} &&
\end{calculs}
\end{correction}

\etape{À vous de jouer}
Mêmes consignes ; pour $I$, $x$ désigne un nombre :
\[
  E=\dfrac{-2}{-3}+\dfrac{-5}{6} \qquad
  F=\dfrac{5}{6}-\dfrac{-3}{4} \qquad
  G=\dfrac{1}{2}-\dfrac{1}{4}+\dfrac{1}{10} \qquad
  H=4-\dfrac{7}{3} \qquad
  I=2-\dfrac{x+3}{4}
\]
\reponse{$E=\frac46-\frac56=-\frac16$ ;\; $F=\frac{10}{12}+\frac{9}{12}=\frac{19}{12}$ ;\;
$G=\frac{10}{20}-\frac{5}{20}+\frac{2}{20}=\frac{7}{20}$ ;\; $H=\frac{12}{3}-\frac73=\frac53$ ;\;
$I=\frac{8-(x+3)}{4}=\frac{5-x}{4}$.}
\end{methodebox}

\afaire{exercices 9 à 12}

\partiecours{Multiplier et diviser des fractions}

\begin{proprietebox}[Propriété 5 --- Produit et quotient]
Pour $b$, $c$ et $d$ non nuls :
\[
  \dfrac{a}{b}\times\dfrac{c}{d}=\dfrac{a\times c}{b\times d}
  \qquad\qquad
  \dfrac{a}{b}\div\dfrac{c}{d}=\dfrac{a}{b}\times\dfrac{d}{c}
\]
Diviser par une fraction, c'est multiplier par son \textbf{inverse}.
\end{proprietebox}

\begin{remarquebox}[Simplifier : seulement des produits]
Dans un quotient de \textbf{produits}, on peut barrer un même facteur en haut et en
bas. Jamais dans une \textbf{somme} : dans $\frac{-2+9}{5+9}$, on ne barre pas les $9$.
Quand on barre tout le numérateur, il reste $1$, et non $0$.
\end{remarquebox}

\begin{methodebox}[Méthode 6 --- Multiplier et diviser des fractions]
On fixe d'abord le signe en comptant les signes moins. On n'effectue pas les produits :
comme pour simplifier une fraction (méthode 4), on décompose chaque nombre en produit de
facteurs premiers, puis on barre les facteurs communs.
\etapeexemples
Calculer et donner le résultat sous forme simplifiée ; pour e), $x$ désigne un nombre :
\begin{questions}[label=\alph*)]
  \item $A=\dfrac{3}{-4}\times\dfrac{-5}{7}$
  \item $B=6\times\dfrac{4}{5}$
  \item $C=\dfrac{28}{15}\times\dfrac{25}{14}$
  \item $D=\dfrac{5}{6}\div\dfrac{-10}{9}$
  \item $E=\dfrac{6x}{5}\times\dfrac{10}{9}$
\end{questions}
\begin{correction}
\cas{a)}
\begin{calculs}
A &= \dfrac{3}{-4}\times\dfrac{-5}{7} &\qquad& \rem{deux signes moins : le résultat est positif}\\
A &= \dfrac{3\times 5}{4\times 7} &\qquad& \rem{numérateurs entre eux, dénominateurs entre eux}\\
A &= \dfrac{15}{28} &\qquad& \rem{aucun facteur commun}
\end{calculs}
\cas{b)}
\begin{calculs}
B &= \dfrac{6}{1}\times\dfrac{4}{5} &\qquad& \rem{un entier est une fraction de dénominateur $1$}\\
B &= \dfrac{24}{5} &&
\end{calculs}
\cas{c)}
\begin{calculs}
C &= \dfrac{28\times 25}{15\times 14} &\qquad& \rem{on n'effectue pas les produits}\\
C &= \dfrac{2\times 2\times 7\times 5\times 5}{3\times 5\times 2\times 7} &\qquad& \rem{$28=2\times 2\times 7$ ; $15=3\times 5$ ; $14=2\times 7$}\\
C &= \dfrac{2\times\barre{2}\times\barre{7}\times\barre{5}\times 5}{3\times\barre{5}\times\barre{2}\times\barre{7}} &\qquad& \rem{on barre un $2$, le $7$ et un $5$}\\
C &= \dfrac{10}{3} &\qquad& \rem{puis on calcule}
\end{calculs}
\cas{d)}
\begin{calculs}
D &= \dfrac{5}{6}\times\dfrac{9}{-10} &\qquad& \rem{on multiplie par l'inverse de $\frac{-10}{9}$}\\
D &= -\dfrac{5\times 3\times 3}{2\times 3\times 2\times 5} &\qquad& \rem{un seul signe moins ; on décompose sans effectuer}\\
D &= -\dfrac{\barre{5}\times\barre{3}\times 3}{2\times\barre{3}\times 2\times\barre{5}} &\qquad& \rem{on barre $5$ et $3$, en haut et en bas}\\
D &= -\dfrac{3}{4} &&
\end{calculs}
\cas{e)}
\begin{calculs}
E &= \dfrac{6x\times 10}{5\times 9} &\qquad& \rem{$6x$ veut dire $6\times x$}\\
E &= \dfrac{2\times 3\times x\times 2\times 5}{5\times 3\times 3} &\qquad& \rem{une lettre est un facteur comme un nombre}\\
E &= \dfrac{2\times\barre{3}\times x\times 2\times\barre{5}}{\barre{5}\times\barre{3}\times 3} &\qquad& \rem{on barre $3$ et $5$}\\
E &= \dfrac{4x}{3} &\qquad& \rem{$2\times x\times 2=4x$}
\end{calculs}
\end{correction}

\etape{À vous de jouer}
Mêmes consignes ; pour $K$, $x$ désigne un nombre :
\[
  F=\dfrac{2}{-7}\times\dfrac{-3}{5} \qquad
  G=4\times\dfrac{2}{9} \qquad
  H=\dfrac{21}{20}\times\dfrac{8}{35}
\]
\[
  I=\dfrac{3}{10}\div\dfrac{-9}{5} \qquad
  J=\dfrac{-4}{9}\div 8 \qquad
  K=\dfrac{4x}{15}\times\dfrac{9}{8}
\]
\reponse{$F=\frac{6}{35}$ ;\; $G=\frac89$ ;\; $H=\frac{6}{25}$ ;\;
$I=\frac{3}{10}\times\frac{5}{-9}=-\frac16$ ;\; $J=\frac{-4}{9}\times\frac18=-\frac{1}{18}$ ;\;
$K=\frac{3x}{10}$.}
\end{methodebox}

\afaire{exercices 13 à 16}

\partiecours{Fractions et priorités de calcul}

\begin{proprietebox}[Propriété 6 --- Priorités de calcul]
On effectue, dans cet ordre :
\begin{enumerate}[label=\arabic*.]
  \item les calculs entre \textbf{parenthèses}, en commençant par les plus intérieures ;
  \item les \textbf{puissances} ;
  \item les \textbf{multiplications et divisions}, de gauche à droite ;
  \item les \textbf{additions et soustractions}, de gauche à droite.
\end{enumerate}
\end{proprietebox}

\begin{proprietebox}[Propriété 7 --- Fraction de fractions]
Pour $b$, $c$ et $d$ non nuls : \quad
$\dfrac{\;\dfrac{a}{b}\;}{\;\dfrac{c}{d}\;}=\dfrac{a}{b}\div\dfrac{c}{d}=\dfrac{a}{b}\times\dfrac{d}{c}$.

\smallskip
Quand le numérateur ou le dénominateur est un calcul, on le calcule à part, puis on
divise.
\end{proprietebox}

\begin{methodebox}[Méthode 7 --- Mener un calcul avec des fractions]
\etapeexemples
Calculer et donner le résultat sous forme simplifiée :
\begin{questions}[label=\alph*)]
  \item $A=\dfrac{5}{6}-\dfrac{-2}{3}\times\dfrac{3}{2^{2}}$
  \item $B=\dfrac{-5}{\;3+\dfrac{3}{2^{2}}\;}$
  \item $C=\left(\dfrac{1}{2}+\dfrac{1}{3}\right)\div\left(\dfrac{3}{4}-\dfrac{1}{6}\right)$
  \item $D=\dfrac{1}{2}+\dfrac{1}{3}\div\dfrac{3}{4}-\dfrac{1}{6}$
\end{questions}
\begin{correction}
\cas{a)}
\begin{calculs}
A &= \dfrac{5}{6}-\dfrac{-2}{3}\times\dfrac{3}{4} &\qquad& \rem{la puissance : $2^{2}=4$}\\
A &= \dfrac{5}{6}-\dfrac{-6}{12} &\qquad& \rem{puis le produit}\\
A &= \dfrac{5}{6}+\dfrac{1}{2} &\qquad& \rem{$\frac{-6}{12}=-\frac12$ ; soustraire un négatif, c'est ajouter}\\
A &= \dfrac{5}{6}+\dfrac{1\fois{3}}{2\fois{3}} &&\\
A &= \dfrac{5}{6}+\dfrac{3}{6} &&\\
A &= \dfrac{8}{6} &&\\
A &= \dfrac{4}{3} &\qquad& \rem{on simplifie par $2$}
\end{calculs}
\cas{b)}
\begin{calculs}
3+\dfrac{3}{2^{2}} &= \dfrac{3\fois{4}}{1\fois{4}}+\dfrac{3}{4} &\qquad& \rem{le dénominateur d'abord, à part}\\
3+\dfrac{3}{2^{2}} &= \dfrac{12}{4}+\dfrac{3}{4} &&\\
3+\dfrac{3}{2^{2}} &= \dfrac{15}{4} &&\\
B &= -5\times\dfrac{4}{15} &\qquad& \rem{diviser par $\frac{15}{4}$, c'est multiplier par $\frac{4}{15}$}\\
B &= -\dfrac{5\times 4}{3\times 5} &&\\
B &= -\dfrac{\barre{5}\times 4}{3\times\barre{5}} &\qquad& \rem{on barre le $5$}\\
B &= -\dfrac{4}{3} &&
\end{calculs}
\cas{c)}
\begin{calculs}
C &= \left(\dfrac{1\fois{3}}{2\fois{3}}+\dfrac{1\fois{2}}{3\fois{2}}\right)\div\left(\dfrac{3\fois{3}}{4\fois{3}}-\dfrac{1\fois{2}}{6\fois{2}}\right) &\qquad& \rem{chaque parenthèse à part}\\
C &= \left(\dfrac{3}{6}+\dfrac{2}{6}\right)\div\left(\dfrac{9}{12}-\dfrac{2}{12}\right) &&\\
C &= \dfrac{5}{6}\div\dfrac{7}{12} &&\\
C &= \dfrac{5}{6}\times\dfrac{12}{7} &\qquad& \rem{on multiplie par l'inverse}\\
C &= \dfrac{5\times 2\times 6}{6\times 7} &&\\
C &= \dfrac{5\times 2\times\barre{6}}{\barre{6}\times 7} &\qquad& \rem{on barre le $6$}\\
C &= \dfrac{10}{7} &&
\end{calculs}
\cas{d)}
\begin{calculs}
D &= \dfrac{1}{2}+\dfrac{1}{3}\times\dfrac{4}{3}-\dfrac{1}{6} &\qquad& \rem{sans parenthèses, la division passe en premier}\\
D &= \dfrac{1}{2}+\dfrac{4}{9}-\dfrac{1}{6} &&\\
D &= \dfrac{1\fois{9}}{2\fois{9}}+\dfrac{4\fois{2}}{9\fois{2}}-\dfrac{1\fois{3}}{6\fois{3}} &\qquad& \rem{dénominateur commun $18$}\\
D &= \dfrac{9}{18}+\dfrac{8}{18}-\dfrac{3}{18} &&\\
D &= \dfrac{14}{18} &&\\
D &= \dfrac{7}{9} &\qquad& \rem{mêmes nombres que $C$, autre résultat}
\end{calculs}
\end{correction}

\etape{À vous de jouer}
Calculer et donner le résultat sous forme simplifiée :
\[
  E=\left(2+\dfrac{1}{3}-\dfrac{3}{4}\right)\times\left(\dfrac{1}{2}-\dfrac{2}{3}\right)
  \qquad
  F=\dfrac{\;\dfrac{1}{3}+\dfrac{1}{4}\;}{\;\dfrac{5}{6}-\dfrac{1}{4}+\dfrac{1}{3}\;}
\]
\[
  G=-\dfrac{3}{7}\times\left[\left(\dfrac{1}{2}+\dfrac{5}{4}\right)-\left(\dfrac{1}{3}-\dfrac{1}{6}\right)\right]
  \qquad
  H=\left(\dfrac{-3}{4}+\dfrac{1}{2}\right)\div\left(\dfrac{2}{3}+\dfrac{5}{-6}\right)
\]
\reponse{$E=\frac{19}{12}\times\left(-\frac16\right)=-\frac{19}{72}$ ;\;
$F=\frac{7}{12}\div\frac{11}{12}=\frac{7}{11}$ ;\;
$G=-\frac37\times\frac{19}{12}=-\frac{19}{28}$ ;\;
$H=\left(-\frac14\right)\div\left(-\frac16\right)=\frac32$.}
\end{methodebox}

\afaire{exercices 17 à 19}

\partiecours{Produit en croix}

\begin{proprietebox}[Propriété 8 --- Produit en croix]
Pour $b\ne 0$ et $d\ne 0$ : \quad si $\dfrac{a}{b}=\dfrac{c}{d}$, alors
$a\times d=b\times c$.
\end{proprietebox}

\begin{remarquebox}[Pourquoi ?]
On écrit les deux fractions avec le dénominateur commun $b\times d$ :
$\dfrac{a\times d}{b\times d}=\dfrac{c\times b}{d\times b}$. Deux fractions égales de même
dénominateur ont le même numérateur : $a\times d=b\times c$.
\end{remarquebox}

\begin{methodebox}[Méthode 8 --- Résoudre une équation par le produit en croix]
\etapeexemples
Résoudre les équations suivantes :
\begin{questions}[label=\alph*)]
  \item $\dfrac{x}{3}=\dfrac{5}{4}$
  \item $\dfrac{x}{4}=\dfrac{x+3}{6}$
  \item $\dfrac{x+3}{2}=\dfrac{2x-1}{3}$
  \item $\dfrac{2x+1}{4}=\dfrac{3x-2}{-2}$
\end{questions}
\begin{correction}
\cas{a)}
\begin{calculs}
4\times x &= 3\times 5 &\qquad& \rem{produit en croix}\\
4x &= 15 &&\\
x &= \dfrac{15}{4} &\qquad& \rem{on divise par $4$ ; la fraction reste telle quelle}
\end{calculs}
\cas{b)}
\begin{calculs}
6x &= 4(x+3) &\qquad& \rem{parenthèses indispensables autour de $x+3$}\\
6x &= 4x+12 &\qquad& \rem{on développe}\\
6x-4x &= 12 &\qquad& \rem{les $x$ à gauche, les nombres à droite}\\
2x &= 12 &&\\
x &= 6 &&
\end{calculs}
\cas{c)}
\begin{calculs}
3(x+3) &= 2(2x-1) &\qquad& \rem{produit en croix}\\
3x+9 &= 4x-2 &\qquad& \rem{on développe les deux membres}\\
9+2 &= 4x-3x &\qquad& \rem{les $x$ du côté où il y en a le plus}\\
x &= 11 &&
\end{calculs}
\cas{d)}
\begin{calculs}
-2(2x+1) &= 4(3x-2) &\qquad& \rem{le $-2$ garde son signe}\\
-4x-2 &= 12x-8 &&\\
-2+8 &= 12x+4x &&\\
6 &= 16x &&\\
x &= \dfrac{6}{16} &&\\
x &= \dfrac{3}{8} &\qquad& \rem{on simplifie par $2$}
\end{calculs}
\end{correction}

\begin{remarquebox}[Vérifier une solution]
On remplace $x$ par la valeur trouvée dans chaque membre. Pour b) :
$\frac{6}{4}=\frac32$ et $\frac{6+3}{6}=\frac96=\frac32$ ; les deux membres sont égaux,
donc $6$ est bien la solution.
\end{remarquebox}

\etape{À vous de jouer}
Résoudre, puis vérifier :
\[
  \text{e) } \dfrac{3}{7}=\dfrac{x}{14} \qquad
  \text{f) } \dfrac{3x}{4}=\dfrac{7-x}{2} \qquad
  \text{g) } \dfrac{2x+1}{3}=\dfrac{x+4}{2} \qquad
  \text{h) } \dfrac{3x-2}{4}=\dfrac{5-x}{8} \qquad
  \text{i) } \dfrac{5x+2}{3}=\dfrac{4-x}{-2}
\]
\reponse{e) $7x=42$, $x=6$ ;\; f)~$6x=28-4x$, $x=\frac{14}{5}$ ;\; g)~$4x+2=3x+12$,
$x=10$ ;\; h)~$24x-16=20-4x$, $x=\frac{36}{28}=\frac97$ ;\; i)~$-10x-4=12-3x$,
$x=-\frac{16}{7}$.}
\end{methodebox}

\afaire{exercices 20 à 22}

\besoinplace{7cm}
\begin{methodebox}[Méthode 9 --- Résoudre une équation à plusieurs fractions]
Quand un membre contient une somme de fractions, le produit en croix ne s'applique
pas : on multiplie \textbf{chaque terme} des deux membres par un dénominateur commun.
\etapeexemples
Résoudre les équations suivantes :
\begin{questions}[label=\alph*)]
  \item $x+\dfrac{x}{3}=12$
  \item $\dfrac{t}{4}-\dfrac{t+1}{6}=1$
  \item $\dfrac{z-1}{2}-\dfrac{z-3}{6}=\dfrac{2z+1}{9}$
\end{questions}
\begin{correction}
\cas{a)}
\begin{calculs}
3\times x+3\times\dfrac{x}{3} &= 3\times 12 &\qquad& \rem{on multiplie chaque terme par $3$, $x$ compris}\\
3x+x &= 36 &&\\
4x &= 36 &&\\
x &= 9 &\qquad& \rem{vérification : $9+\frac93=9+3=12$}
\end{calculs}
\cas{b)}
\begin{calculs}
3t-2(t+1) &= 12 &\qquad& \rem{chaque terme $\times 12$ ; parenthèses à cause du moins}\\
3t-2t-2 &= 12 &&\\
t &= 14 &\qquad& \rem{vérification : $\frac{14}{4}-\frac{15}{6}=3{,}5-2{,}5=1$}
\end{calculs}
\cas{c)}
\begin{calculs}
9(z-1)-3(z-3) &= 2(2z+1) &\qquad& \rem{chaque terme $\times 18$ : $18\div 2=9$, $18\div 6=3$, $18\div 9=2$}\\
9z-9-3z+9 &= 4z+2 &\qquad& \rem{attention : $-3\times(-3)=+9$}\\
6z &= 4z+2 &&\\
2z &= 2 &&\\
z &= 1 &&
\end{calculs}
\end{correction}

\etape{À vous de jouer}
Résoudre :
\[
  \text{d) } \dfrac{x}{3}-\dfrac{x}{5}=4 \qquad
  \text{e) } \dfrac{3t-1}{2}+\dfrac{t-3}{4}=t \qquad
  \text{f) } \dfrac{y}{2}+\dfrac{2y}{3}=\dfrac{y}{4}+11
\]
\[
  \text{g) } \dfrac{y+6}{4}=\dfrac{y-2}{3}-\dfrac{y}{6} \qquad
  \text{h) } \dfrac{3z+1}{4}-\dfrac{5z-2}{3}=\dfrac{1-13z}{12}
\]
\reponse{d) ($\times 15$) $5x-3x=60$, $x=30$ ;\; e) ($\times 4$) $7t-5=4t$, $t=\frac53$ ;\;
f) ($\times 12$) $14y=3y+132$, $y=12$ ;\; g) ($\times 12$) $3y+18=2y-8$, $y=-26$ ;\;
h) ($\times 12$) $-11z+11=1-13z$, $z=-5$.}
\end{methodebox}

\afaire{exercices 23 et 24}

\partiecours{Fractions et problèmes}

\begin{proprietebox}[Propriété 9 --- Fraction d'une quantité]
Prendre les $\dfrac{a}{b}$ \textbf{de} quelque chose, c'est \textbf{multiplier} par
$\dfrac{a}{b}$ : les $\dfrac23$ de $60$ font $\dfrac23\times 60=40$.

\smallskip
Une fraction prise \og du reste\fg{} porte sur ce qui reste, pas sur le total : on
calcule d'abord ce reste.
\end{proprietebox}

\begin{methodebox}[Méthode 10 --- Résoudre un problème de fractions]
\etapeexemples
\begin{questions}[label=\alph*)]
  \item Un train de $360$ places part de Lille avec les trois quarts de ses places
        occupées. À Arras, des voyageurs occupent les deux tiers des places restées
        libres. Combien de places sont encore libres ?
  \item Malo a un forfait mobile de $60$~Go. Il en utilise le quart pour les vidéos,
        puis les deux cinquièmes de ce qui reste pour les réseaux sociaux. Quelle
        fraction du forfait lui reste-t-il ? Combien de gigaoctets cela fait-il ?
  \item Une boutique en ligne vend les deux tiers de son stock de casques audio le
        premier jour des soldes, puis les trois cinquièmes du reste le deuxième jour. Il
        reste alors $20$ casques. Combien en a-t-elle vendu ?
\end{questions}
\begin{correction}
\cas{a)}
\begin{calculs}
\text{occupées à Lille} &= \dfrac34\times 360=270 &\qquad& \rem{les $\frac34$ de $360$}\\
\text{libres} &= 360-270=90 &&\\
\text{occupées à Arras} &= \dfrac23\times 90=60 &\qquad& \rem{les $\frac23$ du reste, pas du total}\\
\text{encore libres} &= 90-60=30 &&
\end{calculs}
\explique{Il reste $30$ places libres.}{on relit la question avant de conclure}
\cas{b)}
\begin{calculs}
\text{reste après les vidéos} &= 1-\dfrac14=\dfrac34 &\qquad& \rem{le forfait entier vaut $1$}\\
\text{réseaux sociaux} &= \dfrac25\times\dfrac34 &\qquad& \rem{les $\frac25$ du reste}\\
\text{réseaux sociaux} &= \dfrac{6}{20}=\dfrac{3}{10} &&\\
\text{reste} &= \dfrac{3\fois{5}}{4\fois{5}}-\dfrac{3\fois{2}}{10\fois{2}} &\qquad& \rem{$\frac34-\frac{3}{10}$ sur le dénominateur $20$}\\
\text{reste} &= \dfrac{15}{20}-\dfrac{6}{20} &&\\
\text{reste} &= \dfrac{9}{20} &&\\
\text{reste en Go} &= \dfrac{9}{20}\times 60=27 &&
\end{calculs}
\explique{Il reste à Malo les $\frac{9}{20}$ du forfait, soit $27$~Go.}{vérification :
$15+18+27=60$}
\cas{c)}
\begin{calculs}
\text{reste après le premier jour} &= 1-\dfrac23=\dfrac13 &\qquad& \rem{on travaille sur des fractions du stock}\\
\text{vendu le deuxième jour} &= \dfrac35\times\dfrac13=\dfrac15 &&\\
\text{reste le soir} &= \dfrac{1\fois{5}}{3\fois{5}}-\dfrac{1\fois{3}}{5\fois{3}} &\qquad& \rem{$\frac13-\frac15$ sur le dénominateur $15$}\\
\text{reste le soir} &= \dfrac{5}{15}-\dfrac{3}{15} &&\\
\text{reste le soir} &= \dfrac{2}{15} &&
\end{calculs}
\explique{Les $\frac{2}{15}$ du stock font $20$ casques, donc $\frac{1}{15}$ en fait
$10$ et le stock est de $15\times 10=150$ casques.}{on remonte au total}
\explique{La boutique en a vendu $150-20=130$.}{la question porte sur les ventes}
\end{correction}

\etape{À vous de jouer}
\begin{questions}[beginpenalty=10000]
  \item Une cuve de $240$~L est remplie aux cinq sixièmes. On en vide les trois quarts.
        Combien de litres reste-t-il dans la cuve ?
  \item Sami dépense le quart de ses économies pour un téléphone, puis les deux tiers de
        ce qui reste pour un vélo. Il lui reste $75$~€. Combien avait-il économisé ?
\end{questions}
\reponse{1\textdegree) $\frac56\times 240=200$~L ; on en vide $\frac34\times 200=150$~L ;
il reste $50$~L.\; 2\textdegree) Reste $\frac34$, vélo $\frac23\times\frac34=\frac12$,
il reste $\frac34-\frac12=\frac14$ des économies : $75\times 4=300$~€.}
\end{methodebox}

\afaire{exercices 25 à 28}

\partiecours{Puissances}

\begin{definitionbox}[Définition 5 --- Puissance d'un nombre]
Pour un nombre $a$ et un entier $n\ge 1$ :
\[
  a^{n}=\underbrace{a\times a\times\dots\times a}_{n\text{ facteurs}}
  \qquad a^{0}=1 \qquad
  a^{-n}=\dfrac{1}{a^{n}} \quad (a\ne 0)
\]
Un exposant positif compte des multiplications, un exposant négatif des divisions :
$5^{-2}=\dfrac{1}{5\times 5}=\dfrac{1}{25}$.

\smallskip
En particulier, $a^{-1}=\dfrac1a$ : $a^{-1}$ est l'\textbf{inverse} de $a$.
\end{definitionbox}

\begin{definitionbox}[Définition 6 --- Puissances de 10]
$10^{n}=1\underbrace{0\dots0}_{n\text{ zéros}}$ \qquad et \qquad
$10^{-n}=\dfrac{1}{10^{n}}=0{,}\underbrace{0\dots0}_{n-1\text{ zéros}}1$
\qquad ($n$ chiffres après la virgule).

\smallskip
Exemples : $10^{6}=1~000~000$ \quad ; \quad $10^{-4}=0{,}0001$.
\end{definitionbox}

\begin{remarquebox}[Signe d'une puissance]
Un nombre négatif élevé à une puissance \textbf{paire} donne un résultat positif ; à
une puissance \textbf{impaire}, un résultat négatif : $(-2)^{4}=16$, $(-2)^{3}=-8$.
Attention aux parenthèses : $(-3)^{2}=9$, mais $-3^{2}=-(3\times 3)=-9$.
\end{remarquebox}

\begin{methodebox}[Méthode 11 --- Calculer avec des puissances]
\etapeexemples
Calculer :
\begin{questions}[label=\alph*)]
  \item $A=(3^{2}-4)\times(-1)^{7}$
  \item $B=(-2)^{3}+(-3)^{2}\times(-2)$
  \item $C=-3^{2}+(-3)^{2}-(-3)^{3}$
\end{questions}
\begin{correction}
\cas{a)}
\begin{calculs}
A &= (9-4)\times(-1) &\qquad& \rem{$3^{2}=9$ ; exposant $7$ impair : $(-1)^{7}=-1$}\\
A &= 5\times(-1) &&\\
A &= -5 &&
\end{calculs}
\cas{b)}
\begin{calculs}
B &= -8+9\times(-2) &\qquad& \rem{$(-2)^{3}=-8$ (impair) ; $(-3)^{2}=9$ (pair)}\\
B &= -8-18 &\qquad& \rem{le produit avant la somme}\\
B &= -26 &&
\end{calculs}
\cas{c)}
\begin{calculs}
C &= -9+9-(-27) &\qquad& \rem{$-3^{2}=-9$ : sans parenthèses, le carré ne porte que sur $3$}\\
C &= -9+9+27 &&\\
C &= 27 &&
\end{calculs}
\end{correction}

\etape{À vous de jouer}
Calculer :
\[
  D=(2^{3}-5)\times(-1)^{6} \qquad E=(-1)^{9}+(-2)^{4}\times(-3) \qquad
  F=(-6)^{2}-6^{2}+(-2)^{3}
\]
\reponse{$D=3\times 1=3$ ;\; $E=-1+16\times(-3)=-49$ ;\; $F=36-36-8=-8$.}
\end{methodebox}

\begin{methodebox}[Méthode 12 --- Passer d'un exposant négatif à une fraction]
\etapeexemples
\begin{questions}[label=\alph*)]
  \item Écrire sous forme de fraction : $A=(-3)^{-2}$ et $B=-4^{-2}$.
  \item Écrire sous la forme $a^{n}$ : $C=\dfrac{1}{9}$, \;
        $D=\dfrac{1}{2\times 2\times 2\times 2\times 2\times 2\times 2}$ et
        $E=\dfrac{1}{(-4)\times(-4)\times(-4)}$.
\end{questions}
\begin{correction}
\cas{a)}
\begin{calculs}
A &= \dfrac{1}{(-3)^{2}} &\qquad& \rem{$a^{-n}=\frac{1}{a^{n}}$}\\
A &= \dfrac{1}{9} &\qquad& \rem{exposant pair : $(-3)^{2}=9$}\\
B &= -\dfrac{1}{4^{2}} &\qquad& \rem{sans parenthèses, l'exposant ne porte que sur $4$}\\
B &= -\dfrac{1}{16} &&
\end{calculs}
\cas{b)}
\begin{calculs}
C &= \dfrac{1}{9^{1}}=9^{-1} &\qquad& \rem{l'inverse de $9$}\\
D &= \dfrac{1}{2^{7}}=2^{-7} &\qquad& \rem{$7$ facteurs égaux à $2$}\\
E &= \dfrac{1}{(-4)^{3}}=(-4)^{-3} &\qquad& \rem{on garde les parenthèses : le moins fait partie du nombre}
\end{calculs}
\end{correction}

\etape{À vous de jouer}
\begin{questions}
  \item Écrire sous forme de fraction : \quad $F=(-2)^{-3}$ \qquad $G=-3^{-2}$ \qquad
        $H=\left(\dfrac{1}{5}\right)^{-1}$
  \item Écrire sous la forme $a^{n}$ : \quad $I=\dfrac{1}{6}$ \qquad
        $J=\dfrac{1}{5\times 5\times 5\times 5}$ \qquad $K=\dfrac{1}{(-7)\times(-7)}$
\end{questions}
\reponse{1\textdegree) $F=\frac{1}{-8}=-\frac18$ ;\; $G=-\frac19$ ;\; $H=5$.\;
2\textdegree) $I=6^{-1}$ ;\; $J=5^{-4}$ ;\; $K=(-7)^{-2}$.}
\end{methodebox}

\afaire{exercices 29 à 34}

\besoinplace{7cm}
\partiecours{Règles de calcul sur les puissances}

\begin{proprietebox}[Propriété 10 --- Règles de calcul]
Pour $a$ et $b$ non nuls, $n$ et $p$ entiers relatifs :
\[
\renewcommand{\arraystretch}{2.2}
\begin{array}{lll}
  a^{n}\times a^{p}=a^{n+p} & \quad & \text{exemple : } 5^{4}\times 5^{-6}=5^{4+(-6)}=5^{-2}\\
  \dfrac{a^{n}}{a^{p}}=a^{n-p} & & \text{exemple : } \dfrac{(-2)^{9}}{(-2)^{4}}=(-2)^{9-4}=(-2)^{5}\\
  \left(a^{n}\right)^{p}=a^{n\times p} & & \text{exemple : } \left(4^{3}\right)^{5}=4^{3\times 5}=4^{15}\\
  (a\times b)^{n}=a^{n}\times b^{n} & & \text{exemple : } (3x)^{2}=3^{2}\times x^{2}=9x^{2}\\
  \left(\dfrac{a}{b}\right)^{n}=\dfrac{a^{n}}{b^{n}} & & \text{exemple : } \left(\dfrac{3}{5}\right)^{2}=\dfrac{3^{2}}{5^{2}}=\dfrac{9}{25}
\end{array}
\]
\end{proprietebox}

\begin{methodebox}[Méthode 13 --- Simplifier avec les règles sur les puissances]
\etapeexemples
Simplifier, puis calculer :
\begin{questions}[label=\alph*)]
  \item $A=\dfrac{(2^{2})^{-3}\times 3}{(2\times 3)^{-2}}$
  \item $B=\dfrac{3^{4}\times 10^{-5}\times(3^{2})^{3}\times(10^{-1})^{2}}{3^{8}\times 10^{-4}}$
\end{questions}
\begin{correction}
\cas{a)}
\begin{calculs}
A &= \dfrac{2^{-6}\times 3^{1}}{2^{-2}\times 3^{-2}} &\qquad& \rem{$(2^{2})^{-3}=2^{-6}$ ; $(2\times 3)^{-2}=2^{-2}\times 3^{-2}$}\\
A &= 2^{-6-(-2)}\times 3^{1-(-2)} &\qquad& \rem{chaque base à part ; on soustrait les exposants}\\
A &= 2^{-4}\times 3^{3} &\qquad& \rem{$-6+2=-4$ et $1+2=3$}\\
A &= \dfrac{27}{16} &\qquad& \rem{$2^{-4}=\frac{1}{16}$}
\end{calculs}
\cas{b)}
\begin{calculs}
B &= \dfrac{3^{4}\times 3^{6}\times 10^{-5}\times 10^{-2}}{3^{8}\times 10^{-4}} &\qquad& \rem{$(3^{2})^{3}=3^{6}$ ; $(10^{-1})^{2}=10^{-2}$}\\
B &= \dfrac{3^{10}\times 10^{-7}}{3^{8}\times 10^{-4}} &\qquad& \rem{on ajoute les exposants des produits}\\
B &= 3^{2}\times 10^{-3} &\qquad& \rem{$10-8=2$ et $-7-(-4)=-3$}\\
B &= 9\times 10^{-3} &\qquad& \rem{soit $0{,}009$}
\end{calculs}
\end{correction}

\etape{À vous de jouer}
Simplifier, puis calculer :
\[
  C=\dfrac{(3^{2})^{-2}\times 2}{(3\times 2)^{-3}} \qquad
  D=\dfrac{2^{5}\times 10^{-4}}{2^{3}\times 10^{-6}} \qquad
  E=\dfrac{a^{4}\times(a^{-3})^{2}}{(a^{2})^{3}} \quad (a\ne 0)
\]
\reponse{$C=3^{-1}\times 2^{4}=\frac{16}{3}$ ;\;
$D=2^{2}\times 10^{2}=400$ ;\; $E=\frac{a^{4}\times a^{-6}}{a^{6}}=a^{-8}$.}
\end{methodebox}

\afaire{exercices 35 à 39}

\partiecours{Notation scientifique}

\begin{definitionbox}[Définition 7 --- Notation scientifique]
La \textbf{notation scientifique} d'un nombre décimal non nul est son écriture
$a\times 10^{n}$, où $a$ a \textbf{un seul chiffre non nul avant la virgule} et $n$
est un entier relatif.

\smallskip
Exemples : $27~800=2{,}78\times 10^{4}$ \quad ; \quad
$-40~500~000=-4{,}05\times 10^{7}$ \quad ; \quad $0{,}000~000~63=6{,}3\times 10^{-7}$.
\end{definitionbox}

\begin{remarquebox}[Pour ne pas se tromper d'exposant]
Le produit $a\times 10^{n}$ ne doit pas changer de valeur : quand la virgule de $a$
recule de $k$ rangs (le nombre devient plus petit), l'exposant augmente de $k$ ; quand
elle avance, l'exposant diminue.
\end{remarquebox}

\begin{methodebox}[Méthode 14 --- Écrire un nombre en notation scientifique]
\etapeexemples
Donner la notation scientifique de :
\begin{questions}[label=\alph*)]
  \item $A=5~600~000$ \quad et \quad $B=0{,}000~47$
  \item $C=318{,}2\times 10^{-6}$
  \item $D=\dfrac{4\times 10^{2}\times 9\times 10^{5}}{60\times 10^{3}}$
  \item $E=6\times 10^{3}+47\times 10^{-2}$
\end{questions}
\begin{correction}
\cas{a)}
\begin{calculs}
A &= 5{,}6\times 10^{6} &\qquad& \rem{la virgule recule de $6$ rangs}\\
B &= 4{,}7\times 10^{-4} &\qquad& \rem{la virgule avance de $4$ rangs}
\end{calculs}
\cas{b)}
\begin{calculs}
C &= 3{,}182\times 10^{2}\times 10^{-6} &\qquad& \rem{$318{,}2=3{,}182\times 10^{2}$}\\
C &= 3{,}182\times 10^{-4} &\qquad& \rem{$2+(-6)=-4$}
\end{calculs}
\cas{c)}
\begin{calculs}
D &= \dfrac{4\times 9}{60}\times\dfrac{10^{2}\times 10^{5}}{10^{3}} &\qquad& \rem{les nombres d'un côté, les puissances de 10 de l'autre}\\
D &= 0{,}6\times 10^{4} &\qquad& \rem{$\frac{36}{60}=0{,}6$ et $2+5-3=4$}\\
D &= 6\times 10^{-1}\times 10^{4} &&\\
D &= 6\times 10^{3} &&
\end{calculs}
\cas{d)}
\begin{calculs}
E &= 6~000+0{,}47 &\qquad& \rem{une somme : on revient à l'écriture décimale}\\
E &= 6~000{,}47 &&\\
E &= 6{,}000~47\times 10^{3} &&
\end{calculs}
\end{correction}

\etape{À vous de jouer}
Donner la notation scientifique de :
\[
  F=3~470~000 \qquad G=0{,}000~902 \qquad H=72{,}6\times 10^{-5} \qquad
  I=\dfrac{8\times 10^{5}\times 3\times 10^{-3}}{6\times 10^{4}}
\]
\reponse{$F=3{,}47\times 10^{6}$ ;\; $G=9{,}02\times 10^{-4}$ ;\;
$H=7{,}26\times 10^{-4}$ ;\; $I=4\times 10^{5-3-4}=4\times 10^{-2}$.}
\end{methodebox}

\afaire{exercices 40 à 44}

\end{document}
